% Comparison figure: Our Framework vs General Medical QA AI
\begin{CJK}{UTF8}{gbsn}
\begin{figure}[H]
\centering
\begin{tikzpicture}[
    node distance=0.7cm,
    box/.style={draw, rounded corners=4pt, minimum width=4.8cm, minimum height=0.65cm, align=center, font=\small},
    header/.style={draw, rounded corners=4pt, fill=black!8, minimum width=4.8cm, minimum height=0.75cm, align=center, font=\small\bfseries},
    yes/.style={draw=green!60!black, fill=green!5, rounded corners=4pt, minimum width=4.8cm, minimum height=0.65cm, align=center, font=\small},
    no/.style={draw=red!50, fill=red!3, rounded corners=4pt, minimum width=4.8cm, minimum height=0.65cm, align=center, font=\small},
    arrow/.style={thick, ->, >=stealth}
]
% Headers
\node[header] (gpt-header) at (0,0) {\large 通用医学问答AI~(e.g., 阿福等)};
\node[header, right=of gpt-header, xshift=1.5cm] (our-header) {\large 我们的框架~(Ours)};

% Feature rows
\node[no, below=0.35cm of gpt-header] (g1) {\textcolor{red!70!black}{\large\textbf{X}}~~单一通用模型};
\node[yes, below=0.35cm of our-header] (o1) {\textcolor{green!60!black}{\large\textbf{\checkmark}}~~每名医生独立优化};

\node[no, below=0.3cm of g1] (g2) {\textcolor{red!70!black}{\large\textbf{X}}~~无科室差异};
\node[yes, below=0.3cm of o1] (o2) {\textcolor{green!60!black}{\large\textbf{\checkmark}}~~按科室/专科适配};

\node[no, below=0.3cm of g2] (g3) {\textcolor{red!70!black}{\large\textbf{X}}~~完全自动化, 无医生参与};
\node[yes, below=0.3cm of o2] (o3) {\textcolor{green!60!black}{\large\textbf{\checkmark}}~~医生历史数据驱动优化};

\node[no, below=0.3cm of g3] (g4) {\textcolor{red!70!black}{\large\textbf{X}}~~单一模型服务所有医生};
\node[yes, below=0.3cm of o3] (o4) {\textcolor{green!60!black}{\large\textbf{\checkmark}}~~多LoRA适配器并行部署};

\node[no, below=0.3cm of g4] (g5) {\textcolor{red!70!black}{\large\textbf{X}}~~无法捕捉个人语言风格};
\node[yes, below=0.3cm of o4] (o5) {\textcolor{green!60!black}{\large\textbf{\checkmark}}~~TIES-MERGING高效融合};

\node[no, below=0.3cm of g5] (g6) {\textcolor{red!70!black}{\large\textbf{X}}~~输出与医生习惯脱节};
\node[yes, below=0.3cm of o5] (o6) {\textcolor{green!60!black}{\large\textbf{\checkmark}}~~个性化响应贴合临床实际};

% Footer
\node[box, draw=red!40, below=0.5cm of g6] (g-footer) {\small 医生需大量修改AI建议};
\node[box, draw=blue!60, fill=blue!5, below=0.5cm of o6] (o-footer) {\small\textbf{医生可直接采用或微调}};

\draw[arrow, red!50, line width=1.5pt] (g-footer.north) -- (g6.south);
\draw[arrow, blue!50, line width=1.5pt] (o-footer.north) -- (o6.south);

\end{tikzpicture}
\caption{我们的个性化框架与通用医学问答AI的对比。通用方案采用单一模型服务所有医生、无科室差异、缺乏医生参与；我们的框架通过医生专属LoRA训练、LAG路由和TIES-MERGING实现每名医生、每个科室的精准个性化优化，输出与医生临床习惯高度一致的响应。}
\label{fig:comparison}
\end{figure}
\end{CJK}